\documentclass[../../main.tex]{subfiles}% 注意这里的文件路径不能用 ./main.tex ,否则用latexmk编译子文件会报错
\graphicspath{{\subfix{./image/}}} % 指定图片目录,后续可以直接使用图片文件名
% 注意这里的文件路径不能用 ../../image/ ,否则用latexmk编译子文件会报错

% 例如:
% \begin{figure}[H]
% \centering
% \includegraphics[scale=0.3]{图.png}
% \caption{}
% \label{figure:图}
% \end{figure}
% 注意:上述\label{}一定要放在\caption{}之后,否则引用图片序号会只会显示??.

\begin{document}

\section{广义函数的运算}

\begin{definition}
设$\Omega\subseteq\mathbb{R}^n$为开集,$A:\mathscr{D}(\Omega)\to\mathscr{D}(\Omega)$为线性算子，若
\begin{align*}
\varphi_j\to\varphi(\mathscr{D}(\Omega))\implies A\varphi_j\to A\varphi(\mathscr{D}(\Omega))\quad(j\to\infty),
\end{align*}
则称$A$是\textbf{连续线性算子}。
\end{definition}

\begin{example}
任意微分算子$\partial^\alpha$是$\mathscr{D}(\Omega)$上的连续线性算子。
\end{example}
\begin{proof}
对任意相对紧集$K\subseteq\Omega$，有$\partial^\alpha:\mathscr{D}_K\to\mathscr{D}_K$，且
\begin{align*}
\|\partial^\alpha\varphi\|_m=\sum\limits_{|\beta|\leqslant m}\max\limits_{x\in K}|\partial^{\beta+\alpha}\varphi(x)|\leqslant\|\varphi\|_{m+|\alpha|}.
\end{align*}
$\partial^\alpha$在$L^2(\Omega)$上不连续，但为闭算子.

\end{proof}

\begin{example}
设$\psi\in C^\infty(\Omega)$，定义乘法算子$A:\varphi\mapsto\psi\cdot\varphi\ (\forall\varphi\in\mathscr{D}(\Omega))$，则$A$是连续线性算子。
\end{example}
\begin{proof}
对任意相对紧集$K\subseteq\Omega$，有$A:\mathscr{D}_K\to\mathscr{D}_K$，且
\begin{align*}
\|\psi\cdot\varphi\|_m&=\sum\limits_{|\alpha|\leqslant m}\max\limits_{x\in K}|\partial^\alpha(\psi\cdot\varphi)(x)|\\
&\leqslant\sum\limits_{|\alpha|\leqslant m}\max\limits_{x\in K}\left|\sum\limits_{\beta\leqslant\alpha}\binom{\alpha}{\beta}\partial^\beta\psi(x)\cdot\partial^{\alpha-\beta}\varphi(x)\right|\\
&\leqslant C(m,\psi)\|\varphi\|_m.
\end{align*}

\end{proof}

\begin{definition}
设$\Omega\subseteq\mathbb{R}^n$为开集,对$\mathscr{D}(\Omega)$上的线性算子$A$，在$\mathscr{D}'(\Omega)$上定义其共轭算子$A^*$满足
\begin{align*}
\langle A^*f,\varphi\rangle=\langle f,A\varphi\rangle\quad(\forall\varphi\in\mathscr{D}(\Omega),\forall f\in\mathscr{D}'(\Omega)).
\end{align*}
\end{definition}

\begin{proposition}
设$\Omega\subseteq\mathbb{R}^n$为开集,则$A^*:\mathscr{D}'(\Omega)\to\mathscr{D}'(\Omega)$为连续线性算子。
\end{proposition}
\begin{proof}
若$f_j\to f\ (\mathscr{D}'(\Omega)),j\to\infty$，则对$\forall\varphi\in\mathscr{D}(\Omega)$，有
\begin{align*}
\langle A^*f_j,\varphi\rangle=\langle f_j,A\varphi\rangle\to\langle f,A\varphi\rangle=\langle A^*f,\varphi\rangle\quad(j\to\infty),
\end{align*}
即$A^*f_j\to A^*f\ (j\to\infty)$.

\end{proof}

\subsection{广义微商}
\begin{definition}
称$\widetilde{\partial}^\alpha=(-1)^{|\alpha|}(\partial^\alpha)^*$为$\alpha$阶\textbf{广义微商运算}，即对$\forall f\in\mathscr{D}'(\Omega)$，
\begin{align*}
\langle\widetilde{\partial}^\alpha f,\varphi\rangle=(-1)^{|\alpha|}\langle f,\partial^\alpha\varphi\rangle\quad(\forall\varphi\in\mathscr{D}(\Omega)).
\end{align*}
\end{definition}

\begin{remark}
若$f(x)\in C^1(\Omega)$，其按自然对应$\langle f,\varphi\rangle=\int_{\Omega}f(x)\varphi(x)\mathrm{d}x\ (\forall\varphi\in\mathscr{D}(\Omega))$生成广义函数。其广义微商$\widetilde{\partial}_{x_i}f$与普通微商$\partial_{x_i}f$对应的广义函数一致，因
\begin{align*}
\langle\widetilde{\partial}_{x_i}f,\varphi\rangle=-\langle f,\partial_{x_i}\varphi\rangle=-\int_{\Omega}f(x)\partial_{x_i}\varphi(x)\mathrm{d}x=\int_{\Omega}\partial_{x_i}f(x)\varphi(x)\mathrm{d}x\quad(\forall\varphi\in\mathscr{D}(\Omega)).
\end{align*}
\end{remark}

\begin{remark}
广义函数对微商运算封闭：任意$f\in\mathscr{D}'(\Omega)$可做任意次广义微商。局部可积函数未必存在普通微商，但作为广义函数总存在广义微商，且广义微商为广义函数，未必是普通函数。
\end{remark}

\begin{remark}
对任意多重指标$\alpha,\beta$，有
\begin{align*}
\widetilde{\partial}^\alpha\cdot\widetilde{\partial}^\beta=\widetilde{\partial}^{\alpha+\beta}=\widetilde{\partial}^\beta\cdot\widetilde{\partial}^\alpha.
\end{align*}
\end{remark}

\begin{remark}
广义微商算子连续，故
\begin{align*}
f_j\to f_0\quad(j\to\infty)\implies\widetilde{\partial}^\alpha f_j\to\widetilde{\partial}^\alpha f_0\quad(j\to\infty),
\end{align*}
即广义微商与极限运算可交换。
\end{remark}

\begin{proposition}
\begin{enumerate}[]
\item 若$Y(x)=\begin{cases}1,&x>0,\\0,&x\leqslant0,\end{cases}$，则$\widetilde{\partial}_x Y(x)=\delta(x)$。

\item 设$\delta^{(\alpha)}$为$\alpha$阶$\delta$广义函数，则$\widetilde{\partial}^\alpha\delta=\delta^{(\alpha)}$。

\item 记$\widetilde{\Delta}=\widetilde{\partial}_{x_1}^2+\cdots+\widetilde{\partial}_{x_n}^2$，则
\begin{align*}
\widetilde{\Delta}|x|^{2-n}&=(2-n)\Omega_n\delta(x)\quad(n\geqslant3),\\
\widetilde{\Delta}\ln|x|&=2\pi\delta(x)\quad(n=2),
\end{align*}
其中$\Omega_n$为$\mathbb{R}^n$中单位球面的面积。
\end{enumerate}
\end{proposition}
\begin{proof}
\begin{enumerate}[(1)]
\item 注意到
\begin{align*}
\langle\widetilde{\partial}_x Y,\varphi\rangle=-\langle Y,\partial_x\varphi\rangle=-\int_{0}^{\infty}\varphi'(x)\mathrm{d}x
=\varphi(0)=\langle\delta,\varphi\rangle\quad(\forall\varphi\in\mathscr{D}(\mathbb{R})).
\end{align*}

\item 注意到
\begin{align*}
\langle\widetilde{\partial}^\alpha\delta,\varphi\rangle=(-1)^{|\alpha|}\langle\delta,\partial^\alpha\varphi\rangle
=(-1)^{|\alpha|}(\partial^\alpha\varphi)(0)=\langle\delta^{(\alpha)},\varphi\rangle.
\end{align*}

\item \begin{enumerate}[(i)]
\item 当$n\geqslant3$时，对$\forall\varphi\in\mathscr{D}(\mathbb{R}^n)$，
\begin{align*}
\langle\widetilde{\Delta}|x|^{2-n},\varphi\rangle&=\langle|x|^{2-n},\Delta\varphi\rangle=\lim\limits_{\varepsilon\to0}\int_{|x|\geqslant\varepsilon}|x|^{2-n}\Delta\varphi(x)\mathrm{d}x\\
&\xlongequal{\text{Green公式}}\lim\limits_{\varepsilon\to0}\left[\int_{|x|\geqslant\varepsilon}\Delta|x|^{2-n}\varphi(x)\mathrm{d}x+\int_{|x|=\varepsilon}\left(\varphi\frac{\partial|x|^{2-n}}{\partial r}-r^{2-n}\frac{\partial\varphi}{\partial r}\right)\mathrm{d}\sigma\right].
\end{align*}
因$\varphi$具有紧支集，可取充分大的球$B(\theta,R)=\{x\in\mathbb{R}^n\mid|x|<R\}$使得$\mathrm{supp}(\varphi)\subseteq B(\theta,R)$。
注意到$\Delta|x|^{2-n}=0\ (|x|\neq0)$，$\varepsilon^{2-n}\int_{|x|=\varepsilon}\frac{\partial\varphi}{\partial r}\mathrm{d}\sigma=O(\varepsilon)\to0\ (\varepsilon\to0)$，且$\varepsilon^{1-n}\int_{|x|=\varepsilon}\varphi(x)\mathrm{d}\sigma\to\varphi(0)\Omega_n\ (\varepsilon\to0)$，故
\begin{align*}
\langle\widetilde{\Delta}|x|^{2-n},\varphi\rangle=(2-n)\varphi(0)\Omega_n=(2-n)\Omega_n\langle\delta,\varphi\rangle\quad(\forall\varphi\in\mathscr{D}(\mathbb{R}^n)).
\end{align*}

\item 当$n=2$时，同理可证
\begin{align*}
\langle\widetilde{\Delta}\ln|x|,\varphi\rangle=2\pi\langle\delta,\varphi\rangle\quad(\forall\varphi\in\mathscr{D}(\mathbb{R}^2)).
\end{align*}
\end{enumerate}
\end{enumerate}
\end{proof}


\subsection{广义函数的乘法}

\begin{definition}
设$\Omega\subseteq\mathbb{R}^n$为开集,对$\forall\psi\in C^\infty(\Omega),\forall f\in\mathscr{D}'(\Omega)$，定义广义函数的乘法满足
\begin{align*}
\langle\psi f,\varphi\rangle=\langle f,\psi\varphi\rangle\quad(\forall\varphi\in\mathscr{D}(\Omega)).
\end{align*}
\end{definition}

\begin{proposition}
设$\Omega\subseteq\mathbb{R}^n$为开集,则广义函数与$C^\infty(\Omega)$函数的乘法算子为连续算子。
\end{proposition}

\begin{example}
\begin{align*}
x^n\widetilde{\partial}^m\delta(x)=
\begin{cases}
(-1)^n\frac{m!}{(m-n)!}\delta^{(m-n)}(x),&m\geqslant n,\\
0,&m<n.
\end{cases}
\end{align*}
\end{example}
\begin{proof}
\begin{align*}
\langle x^n\widetilde{\partial}^m\delta(x),\varphi(x)\rangle&=(-1)^m\langle\delta(x),\partial^m(x^n\varphi(x))\rangle\\
&=(-1)^m\sum\limits_{r=0}^m\binom{m}{r}(\partial^r x^n)(\partial^{m-r}\varphi(x))\bigg|_{x=0}\\
&=
\begin{cases}
\frac{(-1)^m m!n!}{n!(m-n)!}(\partial^{m-n}\varphi)(0),&m\geqslant n,\\
0,&m<n
\end{cases}\\
&=
\begin{cases}
(-1)^n\frac{m!}{(m-n)!}\langle\delta^{(m-n)}(x),\varphi\rangle,&m\geqslant n,\\
0,&m<n.
\end{cases}
\end{align*}

\end{proof}

\begin{remark}
当$f\in L^1_{\mathrm{loc}}(\Omega)$时，$\psi f$的普通函数乘法定义与广义函数乘法定义一致。一般不能定义两个广义函数的乘积，如两个$\delta$函数相乘，结果不再是广义函数。
\end{remark}

\subsection{平移算子与反射算子}
\begin{definition}
对$\forall x_0\in\mathbb{R}^n$，定义\textbf{平移算子}$\tau_{x_0}:\mathscr{D}(\mathbb{R}^n)\to\mathscr{D}(\mathbb{R}^n)$满足
\begin{align*}
(\tau_{x_0}\varphi)(x)=\varphi(x-x_0)\quad(\forall\varphi\in\mathscr{D}(\mathbb{R}^n)).
\end{align*}
\end{definition}

\begin{definition}
定义广义平移算子$\widetilde{\tau}_{x_0}\triangleq(\tau_{-x_0})^*$，即对$\forall f\in\mathscr{D}'(\mathbb{R}^n)$，
\begin{align*}
\langle\widetilde{\tau}_{x_0}f,\varphi\rangle=\langle f,\tau_{-x_0}\varphi\rangle=\langle f,\varphi(x+x_0)\rangle\quad(\forall\varphi\in\mathscr{D}(\mathbb{R}^n)).
\end{align*}
\end{definition}

\begin{remark}
$\widetilde{\tau}_{x_0}$是普通平移算子的推广：若$f(x)\in L^1_{\mathrm{loc}}(\mathbb{R}^n)$，则
\begin{align*}
\int_{\mathbb{R}^n}f(x-x_0)\varphi(x)\mathrm{d}x=\int_{\mathbb{R}^n}f(x)\varphi(x+x_0)\mathrm{d}x=\langle f,\tau_{-x_0}\varphi\rangle=\langle\widetilde{\tau}_{x_0}f,\varphi\rangle\quad(\forall\varphi\in\mathscr{D}(\mathbb{R}^n)),
\end{align*}
即$\widetilde{\tau}_{x_0}f=f(x-x_0)=\tau_{x_0}f$。
\end{remark}

\begin{definition}
定义\textbf{反射算子}$\sigma:\mathscr{D}(\mathbb{R}^n)\to\mathscr{D}(\mathbb{R}^n)$满足
\begin{align*}
(\sigma\varphi)(x)=\varphi(-x)\quad(\forall\varphi\in\mathscr{D}(\mathbb{R}^n)).
\end{align*}
\end{definition}

\begin{definition}
定义广义反射算子$\widetilde{\sigma}\triangleq\sigma^*$，即对$\forall f\in\mathscr{D}'(\mathbb{R}^n)$，
\begin{align*}
\langle\widetilde{\sigma}f,\varphi\rangle=\langle\sigma^*f,\varphi\rangle=\langle f,\sigma\varphi\rangle\quad(\forall\varphi\in\mathscr{D}(\mathbb{R}^n)).
\end{align*}
\end{definition}

\begin{remark}
$\widetilde{\sigma}$是普通反射算子的推广：若$f(x)\in L^1_{\mathrm{loc}}(\mathbb{R}^n)$，则
\begin{align*}
\int_{\mathbb{R}^n}f(-x)\varphi(x)\mathrm{d}x=\int_{\mathbb{R}^n}f(x)\varphi(-x)\mathrm{d}x=\langle f,\sigma\varphi\rangle=\langle\widetilde{\sigma}f,\varphi\rangle\quad(\forall\varphi\in\mathscr{D}(\mathbb{R}^n)),
\end{align*}
即$\widetilde{\sigma}f=f(-x)=\sigma f$。
\end{remark}

\begin{example}
计算:
\begin{enumerate}[(1)]
\item $\widetilde{\partial}_x^n |x|$;
\item $\widetilde{\partial}^n x_+^\lambda(\lambda\in\mathbb R,\lambda\geqslant0)$, 其中
\begin{align*}
x_+^\lambda=
\begin{cases}
x^\lambda, & x>0,\\
0, & x\leqslant0.
\end{cases}
\end{align*}
\end{enumerate}
\end{example}

\begin{example}
求证:
\begin{align*}
\frac{\widetilde{\mathrm d}}{\mathrm d x}\ln|x|=\mathrm{P.V.}\left(\frac{1}{x}\right),
\end{align*}
即
\begin{align*}
\left\langle \frac{\widetilde{\mathrm d}}{\mathrm d x}\ln|x|,\varphi \right\rangle=\lim\limits_{\varepsilon\to0+}\int_{|x|\geqslant\varepsilon}\frac{\varphi(x)}{x}\mathrm{d}x\quad(\forall\varphi\in\mathscr{D}(\mathbb R)).
\end{align*}
\end{example}

\begin{example}
设$\Omega=(\alpha,\beta)\subset\mathbb R$,$x_0\in\Omega$, 又设$f\in C^1(\Omega\setminus\{x_0\})$,$x_0$是$f$的第一类间断点且$f'$在$\Omega\setminus\{x_0\}$内有界. 求证:
\begin{align*}
\frac{\widetilde{\mathrm d}}{\mathrm d x}f=f'+\big(f(x_0+0)-f(x_0-0)\big)\delta(x_0).
\end{align*}
\end{example}

\begin{example}
求证: 对$\forall f\in\mathscr{D}'(\mathbb R^n)$有
\begin{align*}
\widetilde{\partial}_{x_i}f=\lim\limits_{h\to0}\frac{1}{h}\big(\widetilde{\tau}_{-he_i}f-f\big),
\end{align*}
其中
$$e_i=(\underbrace{0,\cdots,0}_{i},1,0,\cdots,0)\quad(i=1,2,\cdots,n).$$
\end{example}

\begin{example}
求证: 对$\forall f\in\mathscr{D}'(\mathbb R^n)$, 以及$\forall\varphi\in\mathscr{D}(\mathbb R^n)$, 函数
\begin{align*}
g(y)=\langle f,\tau_{-y}\varphi\rangle\in C^\infty(\mathbb R^n).
\end{align*}
\end{example}

\begin{example}
求证: 每个$f\in\mathscr{S}'$必是$L^2(\mathbb R^n)$函数乘以多项式的广义微商之有限和, 即$\exists u_\alpha\in L^2(\mathbb R^n)$及偶数$m$, 使得
\begin{align*}
f=(-1)^{|\alpha|}\sum\limits_{|\alpha|\leqslant m}\widetilde{\partial}^\alpha\big[(1+|x|^2)^{\frac{m}{2}}u_\alpha\big].
\end{align*}
提示 利用\eqref{eq:4-2-7}式. \eqref{eq:4-2-7}式中的$m$可以认为是偶数, 这是因为必要的话可在\eqref{eq:4-2-6}式中用$2m$取代$m$.
\end{example}



\end{document}